\begin{table}[ht]
\centering
\caption{Predicted probability of usable height data by wealth quintile and survey period}
\label{tab:pp_wealth}
\begin{threeparttable}

\small
\renewcommand{\arraystretch}{1.15}
\setlength{\tabcolsep}{10pt}

\begin{tabularx}{\textwidth}{>{\raggedright\arraybackslash}X r r}
\toprule

\textbf{Wealth quintile} & \textbf{2011--2016} & \textbf{2017--2022} \\
\midrule
Poorest & 0.918 (0.915, 0.922) & 0.925 (0.922, 0.929) \\
Poorer  & 0.921 (0.918, 0.925) & 0.930 (0.926, 0.934) \\
Middle  & 0.917 (0.912, 0.921) & 0.929 (0.925, 0.933) \\
Richer  & 0.915 (0.909, 0.922) & 0.927 (0.922, 0.931) \\
Richest & 0.903 (0.897, 0.909) & 0.902 (0.896, 0.908) \\
\bottomrule
\end{tabularx}

\begin{tablenotes}[flushleft]
\footnotesize
\item Notes: Predicted probabilities are marginal estimates from a survey-weighted logistic regression including a wealth quintile $\times$ survey period interaction. Values are predicted probabilities with 95\% confidence intervals. Joint Wald test for interaction: \(p=0.0487\).
\end{tablenotes}

\end{threeparttable}
\end{table}